
\documentclass[letterpaper,11pt]{article}
%%%%%%%%%%%%%%%%%%%%%%%%%%%%%%%%%%%%%%%%%%%%%%%%%%%%%%%%%%%%%%%%%%%%%%%%%%%%%%%%%%%%%%%%%%%%%%%%%%%%%%%%%%%%%%%%%%%%%%%%%%%%%%%%%%%%%%%%%%%%%%%%%%%%%%%%%%%%%%%%%%%%%%%%%%%%%%%%%%%%%%%%%%%%%%%%%%%%%%%%%%%%%%%%%%%%%%%%%%%%%%%%%%%%%%%%%%%%%%%%%%%%%%%%%%%%
\usepackage{graphicx}
\usepackage{amsmath,amsfonts,amssymb,amsthm,float}
\usepackage{hyperref}
\usepackage[utf8]{inputenc}
\usepackage[left=2cm, right=2cm, top=2cm, bottom=2cm]{geometry}

\setcounter{MaxMatrixCols}{10}
%TCIDATA{OutputFilter=LATEX.DLL}
%TCIDATA{Version=5.50.0.2953}
%TCIDATA{<META NAME="SaveForMode" CONTENT="1">}
%TCIDATA{BibliographyScheme=BibTeX}
%TCIDATA{LastRevised=Saturday, November 21, 2026 07:06:02}
%TCIDATA{<META NAME="GraphicsSave" CONTENT="32">}
%TCIDATA{ComputeDefs=
%$\alpha =0.5449$
%$\beta =2.761\times 10^{-6}$
%$\delta =2.6486\times 10^{-6}$
%$\gamma =0.8464$
%}


\input{tcilatex}
\renewcommand{\baselinestretch}{1.15}
\setlength{\parindent}{0pt}
\setlength{\parskip}{0.5\baselineskip}
\pretolerance=2000 \tolerance=3000
\renewcommand{\abstractname}{Resumen}

\begin{document}

\title{Pr\'{a}ctica 2: Caos en Sistemas Biol\'{o}gicos}
\author{Valadez Madera Nestor Roman $M26210087$ \\
%EndAName
Departamento de Ingenier\'{\i}a El\'{e}ctrica y Electr\'{o}nica\\
Tecnol\'{o}gico Nacional de M\'{e}xico / Instituto Tecnol\'{o}gico de Tijuana%
}
\maketitle

\noindent \textbf{Palabras clave: }Modelo Matematico; Ecuaciones
diferenciales; Simulaciones numericas; Lotka-Volterra; MATLAB.

\bigskip

\noindent Correo: \textbf{m26210087@tectijuana.edu.mx}

\noindent \noindent Carrera: \textbf{Maestr\'{\i}a en\ Ciencias de la
Ingenier\'{\i}a.}

\noindent Asignatura: \textbf{Modelado Matem\'{a}tico}

\noindent Profesor: \href{https://biomath.xyz/}{\textbf{Dr. Paul Antonio
Valle Trujillo}} (paul.vt@tijuana.tecnm.mx)

\bigskip

\section{Modelo matematico}

\qquad El modelo matem\'{a}atico que describe la din\'{a}mica entre tres
poblaciones celulares est\'{a} descrito por las siguientes tres EDOs no
lineales de primer orden:

\begin{eqnarray*}
\dot{x} &=&r_{1}x(1-b_{1}x)-a_{12}xy-a_{13}xz, \\
\dot{y} &=&r_{2}y(1-b_{2}y)-a_{21}xy, \\
\dot{z} &=&(r_{3}-a_{31})xz-d_{3}z+\rho ,
\end{eqnarray*}

donde todos los par\'{a}metros son positivos, las condiciones iniciales no
negativas y el par\'{a}metro de control $\rho >0$.

\bigskip

\section{Positividad del sistema}

Con base en el trabajo de De Leenherr y Aeyels, se realiza el siguiente c%
\'{a}lculo para establecer la positividad del sistema:

\begin{eqnarray*}
\left. \dot{x}\right\vert _{x=0}
&=&r_{1}(0)(1-b_{1}(0))-a_{12}(0)y-a_{13}(0)z=0, \\
\left. \dot{y}\right\vert _{y=0} &=&r_{2}(0)(1-b_{2}(0))-a_{21}x(0)=0, \\
\left. \dot{z}\right\vert _{z=0} &=&(r_{3}-a_{31})x(0)-d_{3}(0)+\rho =\rho
\geq 0,
\end{eqnarray*}%
entonces, siguiendo el criterio general para la positividad en las
soluciones de un sistema din\'{a}mico no lineal, se establece el siguiente
resultado:

\QTP{Body Math}
\textbf{Resultado. Positividad e invariancia en el octante no negativo. }%
\textit{Sea }$\varphi =(x(t),y(t),z(t))^{T}$\textit{\ una soluci\'{o}n del
sistema con condiciones iniciales no negativas, es decir, }$\varphi
(t_{0})\in 
%TCIMACRO{\U{211d} }%
%BeginExpansion
\mathbb{R}
%EndExpansion
_{+,0}^{3}$\textit{, por lo tanto, el octante no negativo definido por%
\newline
}%
\begin{equation*}
%TCIMACRO{\U{211d} }%
%BeginExpansion
\mathbb{R}
%EndExpansion
_{+,0}^{3}=\{\varphi \in 
%TCIMACRO{\U{211d} }%
%BeginExpansion
\mathbb{R}
%EndExpansion
_{+,0}^{3}:x,y,z\geq 0\},
\end{equation*}

\QTP{Body Math}
\textit{es positivamente invariante y }$\varphi (t,x,y,z)\in 
%TCIMACRO{\U{211d} }%
%BeginExpansion
\mathbb{R}
%EndExpansion
_{+,0}^{3}$\textit{\ para }$t\geq 0.$

\QTP{Body Math}
$\bigskip $

\section{Puntos de equilibrio libre de c\'{e}lulas patol\'{o}gicas}

Con el objetivo de determinar condiciones para la eliminaci\'{o}n de c\'{e}%
lulas patol\'{o}gicas, se calculan los equilibrios donde $x=0$, es decir, se
debe resolver el siguiente sistema de ecuaciones:

\begin{eqnarray*}
r_{1}x(1-b_{1}x)-a_{12}xy-a_{13}xz &=&0, \\
r_{2}y(1-b_{2}y)-a_{21}xy &=&0, \\
(r_{3}-a_{31})xa-d_{3}z+\rho &=&0,
\end{eqnarray*}%
entonces, si $x=0$, se reescribe de la siguiente forma:

\begin{eqnarray*}
x &=&0 \\
r_{2}y(1-b_{2}y) &=&0 \\
-d_{3}z+\rho &=&0
\end{eqnarray*}%
por lo tanto, se resuelve despejando cada variable de estado:

\begin{eqnarray*}
x^{\ast } &=&0, \\
y(1-b_{2}y) &=&0, \\
z^{\ast } &=&\frac{\rho }{d_{3}},
\end{eqnarray*}%
resolviendo para $y$ se obtiene lo siguiente:

\begin{eqnarray*}
y_{0}^{\ast } &=&0, \\
y_{1}^{\ast } &=&\frac{1}{b_{2}},
\end{eqnarray*}%
y los equilibrios donde $x=0$, est\'{a}an dados por:

\begin{eqnarray*}
(x_{0}^{\ast },y_{0}^{\ast },z_{0}^{\ast }) &=&(0,0,\frac{\rho }{d_{3}}), \\
(x_{1}^{\ast },y_{1}^{\ast },z_{1}^{\ast }) &=&(0,\frac{1}{b_{2}},\frac{\rho 
}{d_{3}}),
\end{eqnarray*}

\section{Estabilidad asint\'{o}tica local}

Ahora, con el proposito de establecer condiciones para asegurar la
estabilidad asint\'{o}tica local del punto de equilibrio libre de c\'{e}%
lulas patol\'{o}gicas $(x_{1}^{\ast },y_{1}^{\ast },z_{1}^{\ast })$, se
calcula la matriz Jacobiana del sistema como se muestra a continuaci\'{o}n:

\begin{equation*}
J=\left[ 
\begin{array}{ccc}
\frac{\partial \dot{x}}{\partial x} & \frac{\partial \dot{x}}{\partial y} & 
\frac{\partial \dot{x}}{\partial z} \\ 
\frac{\partial \dot{y}}{\partial x} & \frac{\partial \dot{y}}{\partial y} & 
\frac{\partial \dot{y}}{\partial z} \\ 
\frac{\partial \dot{z}}{\partial x} & \frac{\partial \dot{z}}{\partial y} & 
\frac{\partial \dot{z}}{\partial z}%
\end{array}%
\right] =\left[ 
\begin{array}{ccc}
r_{1}-a_{12}y-a_{13}z-2b_{1}r_{1}x & -a_{12}x & -a_{13}x \\ 
-a_{21}y & r_{2}-a_{21}x-2b_{2}r_{2}y & 0 \\ 
r_{3}z-a_{31}z & 0 & r_{3}x-a_{31}x-d_{3}%
\end{array}%
\allowbreak \right]
\end{equation*}%
Para evaluar la estabilidad local del equilibrio de inter\'{e}s, se eval\'{u}%
a $J$ en $(x_{1}^{\ast },y_{1}^{\ast },z_{1}^{\ast })$ y se calculan los
valores propios de la matriz resultante, es decir,

\begin{eqnarray*}
\left. J\right\vert _{(x_{1}^{\ast },y_{1}^{\ast },z_{1}^{\ast })} &=&\left[ 
\begin{array}{ccc}
r_{1}-a_{12}y-a_{13}z-2b_{1}r_{1}x & -a_{12}x & -a_{13}x \\ 
-a_{21}y & r_{2}-a_{21}x-2b_{2}r_{2}y & 0 \\ 
r_{3}z-a_{31}z & 0 & r_{3}x-a_{31}x-d_{3}%
\end{array}%
\right] _{x=0,y=1/b_{2},z=\rho /d_{3}} \\
\left. J\right\vert _{(x_{1}^{\ast },y_{1}^{\ast },z_{1}^{\ast })} &=&\left[ 
\begin{array}{ccc}
r_{1}-\frac{a_{12}}{b_{2}}-\frac{a_{13}\rho }{d_{3}} & 0 & 0 \\ 
-\frac{a_{21}}{b_{2}} & -r_{2} & 0 \\ 
\frac{\rho }{d_{3}}(r_{3}-a_{31}) & 0 & -d_{3}%
\end{array}%
\right]
\end{eqnarray*}%
Para determinar los valores propios de la matriz, se c\'{a}lculan las ra%
\'{\i}ces del Jacobiano:

\begin{eqnarray*}
\det (A-\lambda I) &=&0, \\
\det (%
\begin{array}{ccc}
r_{1}-\frac{a_{12}}{b_{2}}-\frac{a_{13}\rho }{d_{3}} & 0 & 0 \\ 
-\frac{a_{21}}{b_{2}} & -r_{2} & 0 \\ 
\frac{\rho }{d_{3}}(r_{3}-a_{31}) & 0 & -d_{3}%
\end{array}%
-%
\begin{array}{ccc}
\lambda & 0 & 0 \\ 
0 & \lambda & 0 \\ 
0 & 0 & \lambda%
\end{array}%
) &=&0, \\
-\frac{1}{b_{2}d_{3}}\left( \lambda +d_{3}\right) \left( \lambda
+r_{2}\right) \left( d_{3}a_{12}+\lambda b_{2}d_{3}+\rho
b_{2}a_{13}-b_{2}d_{3}r_{1}\right) \allowbreak &=&0 \\
\left( \lambda +d_{3}\right) \left( \lambda +r_{2}\right) \left(
d_{3}a_{12}+\lambda b_{2}d_{3}+\rho b_{2}a_{13}-b_{2}d_{3}r_{1}\right) &=&0
\end{eqnarray*}%
Cuando la matriz $A$ es triangular (inferior o superior), entonces, los
valores propios son los elementos de la diagonal:

\begin{eqnarray*}
\lambda _{1} &=&r_{1}-\frac{a_{12}}{b_{2}}-\frac{a_{13}\rho }{d_{3}}, \\
\lambda _{2} &=&-r_{2}, \\
\lambda _{3} &=&-d_{3},
\end{eqnarray*}%
Se observa que $\lambda _{2,3}<0$ y se determina la siguiente condici\'{o}n
sobre el par\'{a}metro de control $\rho $ para asegurar que $\lambda _{1}<0$:

\begin{eqnarray*}
r_{1}-\frac{a_{12}}{b_{2}}-\frac{a_{13}\rho }{d_{3}} &<&0, \\
\rho &>&\frac{d_{3}}{a_{13}}(r_{1}-\frac{a_{12}}{b_{2}}),
\end{eqnarray*}%
por lo tanto, se establece el siguiente resultado con respecto a la
estabilidad asint\'{o}tica local de $(x_{1}^{\ast },y_{1}^{\ast
},z_{1}^{\ast })$

\bigskip

\QTP{Body Math}
\textbf{Resultado: Condici\'{o}n de eliminaci\'{o}n local de c\'{e}lulas
patol\'{o}gicas.} Si la condici\'{o}n

\QTP{Body Math}
\begin{equation*}
\rho >\frac{d_{3}}{a_{13}}(r_{1}-\frac{a_{12}}{b_{2}}),
\end{equation*}%
se cumple, entonces, el equilibrio libre de c\'{e}lulas patol\'{o}gicas $%
(x_{1}^{\ast },y_{1}^{\ast },z_{1}^{\ast })$ es local asint\'{o}ticamente
estable y se puede asegurar la eliminaci\'{o}n de una poblaci\'{o}n inicial
de c\'{e}lulas patol\'{o}gicas cercana a dicho equilibrio.

\bigskip

Del resultado anterior se ontiene el siguiente resultado complementario:

\bigskip

\QTP{Body Math}
\textbf{Resultado: Persistencia. }\textit{Si la siguiente condici\'{o}n se
cumple:}

\begin{equation*}
\rho <\frac{d_{3}}{a_{13}}(r_{1}-\frac{a_{12}}{b_{2}}),
\end{equation*}

\QTP{Body Math}
entonces, se observar\'{a} una persistencia de la poblaci\'{o}n de c\'{e}%
lulas patol\'{o}gicas en las soluciones y trayectorias del sistema.

\QTP{Body Math}
$\bigskip $

Tarea: Realizar el an\'{a}lisis de estabilidad asint\'{o}tica local de $%
(x_{0}^{\ast },y_{0}^{\ast },z_{0}^{\ast })=(0,0,\frac{\rho }{d_{3}})$

\bigskip

\begin{eqnarray*}
\left. J\right\vert _{(x_{1}^{\ast },y_{1}^{\ast },z_{1}^{\ast })} &=&\left[ 
\begin{array}{ccc}
r_{1}-a_{12}y-a_{13}z-2b_{1}r_{1}x & -a_{12}x & -a_{13}x \\ 
-a_{21}y & r_{2}-a_{21}x-2b_{2}r_{2}y & 0 \\ 
r_{3}z-a_{31}z & 0 & r_{3}x-a_{31}x-d_{3}%
\end{array}%
\right] _{x=0,y=0,z=\rho /d_{3}} \\
\left. J\right\vert _{(x_{1}^{\ast },y_{1}^{\ast },z_{1}^{\ast })} &=&\left[ 
\begin{array}{ccc}
r_{1}-\frac{a_{13}\rho }{d_{3}} & 0 & 0 \\ 
0 & r_{2} & 0 \\ 
\frac{r_{3}\rho }{d_{3}}-\frac{a_{31}\rho }{d_{3}} & 0 & -d_{3}%
\end{array}%
\right]
\end{eqnarray*}

Para determinar los valores propios de la matriz, se c\'{a}lculan las ra%
\'{\i}ces del Jacobiano:

\begin{eqnarray*}
\det (A-\lambda I) &=&0, \\
\det (%
\begin{array}{ccc}
r_{1}-\frac{a_{13}\rho }{d_{3}} & 0 & 0 \\ 
0 & r_{2} & 0 \\ 
\frac{r_{3}\rho }{d_{3}}-\frac{a_{31}\rho }{d_{3}} & 0 & -d_{3}%
\end{array}%
-%
\begin{array}{ccc}
\lambda & 0 & 0 \\ 
0 & \lambda & 0 \\ 
0 & 0 & \lambda%
\end{array}%
) &=&0,
\end{eqnarray*}

\begin{eqnarray*}
\lambda _{1} &=&r_{1}-\frac{a_{13}\rho }{d_{3}}, \\
\lambda _{2} &=&r_{2}, \\
\lambda _{3} &=&-d_{3},
\end{eqnarray*}

\bigskip Se observa que $\lambda _{1}<0$ si se cumple la condicion $\rho >%
\frac{r_{1}d_{3}}{a_{13}}$, $\lambda _{2}>0$ y $\lambda _{3}<0$. El m\'{e}%
todo de estabilidad de Lyapunov nos indica que para $\lambda >0$ la soluci%
\'{o}n no converge al punto de equilibrio y se presentar\'{a} una dinamica
del tipo Saddle.

\section{Conjuntos Compactos Invariantes}

\qquad Para definir el dominio de loscalizaci\'{o}n del sistema bajo
estudio, es necesario proponer diferentes funciones localizadoras hasta
encontrar las que permitan establecer l\'{\i}mites inferiores y/o superiores
de cada una de las variables de estado, en este caso, $x(t)$, $y(t)$ y $z(t)$%
. Para sistemas biol\'{o}gicos cuya din\'{a}mica se localiza en el ortante
no negativo, se debe recordar que las soluciones con condiciones iniciales
no negativas se encontrar\'{a}n en el siguiente dominio:%
\begin{equation*}
%TCIMACRO{\U{211d} }%
%BeginExpansion
\mathbb{R}
%EndExpansion
_{+,0}^{3}=\{\varphi \in 
%TCIMACRO{\U{211d} }%
%BeginExpansion
\mathbb{R}
%EndExpansion
_{+,0}^{3}:x,y,z\geq 0\}.
\end{equation*}

Primero, se explora la funci\'{o}n

\begin{equation*}
h_{1}=x,
\end{equation*}

y se calcula su derivada de Lie (la derivada parcial con respecto al tiempo)

\begin{equation*}
L_{f}h_{1}=\frac{dx}{dt}=r_{1}x(1-b_{1}x)-a_{12}xy-a_{13}xz,
\end{equation*}

a partir de lo anterior, se define el conjunto $S(h_{1})=\{L_{f}h_{1}=0\}$,
por lo tanto,%
\begin{equation*}
S(h_{1})=\{r_{1}x(1-b_{1}x)-a_{12}xy-a_{13}xz=0\},
\end{equation*}

se observa que la bvariable de estado $x(t)$ se encuentra como factor com%
\'{u}n, por lo tanto, se reescribe el conjunto de la siguiente forma

\begin{equation*}
S(h_{1})=\{r_{1}-b_{1}r_{1}x-a_{12}y-a_{13}z=0\}\cup \{x=0\},
\end{equation*}

del resultado anterior, es evidente que el l\'{\i}mite inferior de la
variable est\'{a} dado por:

\begin{equation*}
x(t)\geq x_{\inf }=0,
\end{equation*}

y adem\'{a}s, se establece como plano invariante o semipermeable, por que
las trayectorias/soluciones no pueden atravesar el plano $x=0$ desde valores
positivos hacia los negativos. Ahora, para encontrar un l\'{\i}mite
superior, es necesario verificar si es posible reescribir el conjunto de la
siguiente forma general:

\QTP{Body Math}
\begin{equation*}
S(h)=\{x=\text{constante}-f(x,y,z)\}\cup \{x=0\},
\end{equation*}

con base en lo anterior, se reescribe $S(h_{l})$ de la siguiene forma

\begin{equation*}
S(h_{1})=\{x=\frac{1}{b_{1}}-\frac{a_{12}}{b_{1}r_{1}}y-\frac{a_{13}}{%
b_{1}r_{1}}z\}\cup \{x=0\},
\end{equation*}

donde

\begin{equation*}
f(y,z)=\frac{a_{12}}{b_{1}r_{1}}y+\frac{a_{13}}{b_{1}r_{1}}z
\end{equation*}

por lo tanto

\begin{equation*}
S(h_{1})=\{x=\frac{1}{b_{1}}-f(y,z)\}\cup \{x=0\},
\end{equation*}

con base en lo anterior, se concluye el siguiente resultado para los l\'{\i}%
mites de la variable $x(t):$

\begin{equation*}
K_{x}=\{x_{\inf }=0\leq x(t)\leq x_{\sup }=\frac{1}{b_{1}}\}.
\end{equation*}

\qquad Ahora se muestra el an\'{a}lisis para determinar los l\'{\i}mites de
la variable $y(t):$

\QTP{Body Math}
\begin{eqnarray*}
h_{2} &=&y, \\
L_{f}h_{2} &=&\frac{dy}{dt}=r_{2}y(1-b_{2}y)-a_{21}xy, \\
S(h_{2}) &=&\{L_{f}h_{2}=0\}, \\
S(h_{2} &=&\{r_{2}y(1-b_{2}y)-a_{21}xy=0\} \\
S(h_{2}) &=&\{r_{2}-b_{2}r_{2}y-a_{21}x=0\}\cup \{y=0\}, \\
S(h_{2}) &=&\{y=\text{constante}-f(x,y,z)\}\cup \{y=0\}, \\
S(h_{2}) &=&\{y=\frac{1}{b_{2}}-\frac{a_{21}}{b_{2}r_{2}}x\}\cup \{y=0\}, \\
f(x) &=&\frac{a_{21}}{b_{2}r_{2}}x, \\
S(h_{2}) &=&\{y=\frac{1}{b_{2}}-f(x)\}\cup \{y=0\}, \\
K_{y} &=&\{y_{\inf }=0\leq y(t)\leq y_{\sup }=\frac{1}{b_{2}}\}.
\end{eqnarray*}

$\qquad $Para determinar el l\'{\i}mite inferior de la variable de estado $%
z(t)$, se propone la siguiente funci\'{o}n localizadora:

\begin{eqnarray*}
h_{3} &=&z, \\
L_{f}h_{3} &=&\frac{dz}{dt}=(r_{3}-a_{31})xz-d_{3}z+\rho , \\
S(h_{3}) &=&\{L_{f}h_{3}=0\}, \\
S(h_{3}) &=&\{(r_{3}-a_{31})xz-d_{3}z+\rho ,
\end{eqnarray*}

esta funci\'{o}n se propone para encontrar el l\'{\i}mite inferior, por lo
tanto, la forma general, se define como se muestra a continuaci\'{o}n:

\QTP{Body Math}
\begin{equation*}
S(h)=\{z=\text{constante}+f(x,y,z)\}
\end{equation*}

\QTP{Body Math}
\begin{equation*}
S(h_{3})=\{z=\frac{\rho }{d_{3}}+\frac{r_{3}-a_{31}}{d_{3}}xz\},
\end{equation*}

entonces, si la siguiente condicion se cumple

\begin{equation*}
r_{3}-a_{31}>0,
\end{equation*}

se reescribe el conjunto como indica

\begin{equation*}
S(h_{3})=\{z=\frac{\rho }{d_{3}}+f(x,z)\},
\end{equation*}

y se define el siguiente resultado

\begin{equation*}
K_{z}=\{z(t)\geq z_{\inf }=\frac{\rho }{d_{3}}\}.
\end{equation*}

\bigskip

\section{Estailidad asint\'{o}tica global: Condiciones de eliminaci\'{o}n}

\qquad Para establecer condiciones de eliminaci\'{o}n de la variable que
describe la din\'{a}mica de la poblaci\'{o}n de c\'{e}ulas patol\'{o}gicas,
se debe apicar el m\'{e}todo directo de Lyapunov, por lo cual, se debe
proponer la denominada funci\'{o}n candidata de Lyapunov:

\begin{equation*}
V=x,
\end{equation*}

\qquad Nota: Se toma la variable de las c\'{e}lulas patol\'{o}gicas porque
es la din\'{a}mica que se quiere controlar, es decir, enviar al origen $x=0$%
, esto implica que se elimina completamente la poblaci\'{o}n de c\'{e}lulas
patol\'{o}gicas en el sistema.

\qquad Se calcula la derivada teporal de la funci\'{o}n candidata de
Lyapunov como se indica:

\begin{equation*}
\dot{V}=\frac{dV}{dt}=\dot{x}=r_{1}x(1-b_{1}x)-a_{12}xy-a_{13}xz,
\end{equation*}

para cumplir con las condiciones del m\'{e}todo directo de Lyapunov, es
decir,

1. La funci\'{o}n candidata debe ser definida postiva: Que significa
definida positiva?

Significa que la funcion candidata $V(x)>0$ para $x>0$ y $x\neq 0$

2. La funcion candidata debe de ser radialmente desacotada: Que significa
radialmente desacotada?

Significa que $V(x)\longrightarrow \infty $ cuando $\left\Vert x\right\Vert
\longrightarrow \infty $

3. La derivada temporal de la funcion candidata debe de ser al menos
semidefinida negativa, preferentemente negativa definida. Que significa
semidefinida negativa y definida negativa?

Es semidefinida negativa si $\dot{V}(0)=0$ y $\dot{V}\leq 0$ para $x>0$. Es
definida negativa si $\dot{V}<0$ para $x>0$ y $\dot{V}(0)=0$.

4. Que sucede cuando la derivada es semidefinida negativa? Que metodo se
debe aplicar para establecer las condiciones de estabilidad global?

Cuando la derivada temporal de una funci\'{o}n candidata de Lyapunv es
semidefinida negativa, se garantiza que la funci\'{o}n no aumenta a lo largo
de las trayectorias del sistema. Para establecer condiciones de estabilidad
asint\'{o}tica global se puede aplicar el principio de invariancia de
LaSalle, el cual permite \qquad \qquad analizar el comportamiento de las
trayectorias dentro del conjunto deonde la derivada temporal de la funci\'{o}%
n candidata es cero.

Con base en el resultado de la derivada, se evalua en los l\'{\i}mites
encontrados para cada variale:

\begin{equation*}
\left. \dot{V}\right\vert _{K_{x}\cap K_{y}\cap K_{z}},
\end{equation*}

obetniendo el siguiente resultado

\begin{equation*}
\dot{V}=(r_{1}-b_{1}r_{1}x-a_{12}y-a_{13}z)x\leq r_{1}-b_{1}r_{1}x_{\inf
}-a_{12}y_{\inf }-a_{13}z_{\inf })x\leq 0,
\end{equation*}

de la condicion anterior, se formula la siguiente desigualdad

\begin{equation*}
r_{1}-\frac{a_{13}}{d_{3}}\rho <0,
\end{equation*}

se resuelve la desigualdad con respecto al par\'{a}metro $\rho $

\begin{equation*}
\rho >\frac{r_{1}d_{3}}{a_{13}}.
\end{equation*}

\qquad Con base en la condici\'{o}n anterior, se concluye el siguiente
resultado:

\bigskip 

\bigskip \textbf{Resultado: Condici\'{o}n de eliminaci\'{o}n.} \textit{Si la
siguiente condici\'{o}n se satisface }

\begin{equation*}
\rho >\frac{r_{1}d_{3}}{a_{13}}
\end{equation*}

entonces, la poblaci\'{o}n de c\'{e}lulas patol\'{o}gicas se eliminar\'{a}
para cualquier condicion inicial positiva, adem\'{a}as, las trayectorias
converger\'{a}n al siguiente punto de equilibrio:

\begin{equation*}
(x_{1}^{\ast },y_{1}^{\ast },z_{1}^{\ast })=(0,\frac{1}{b_{2}},\frac{\rho }{%
d_{3}}).
\end{equation*}

\end{document}
